\documentclass{article}

\usepackage{amsmath,amssymb}
\usepackage{xurl}
\usepackage[hidelinks]{hyperref}
\setlength{\emergencystretch}{2em}

\input{./command}

\title{The Perron--Frobenius Rank-One Step in Cirac--Perez-Garcia--Schuch--Verstraete 2017 Appendix C.2}
\author{}
\date{2026-04-30}

\begin{document}
\maketitle
\gapnote{false-source}{open}

\section{The assertion}
\label{sec:cpgsv17_assertion}

Appendix C.2 of Cirac--P\'erez-Garc\'ia--Schuch--Verstraete relates the strong
area law and zero correlation length for a simple matrix product density
operator to an outer-product factorization of a finite trace matrix
\cite[Appendix~C.2, Lemmas~C.3--C.4]{Cirac2016MPDO_arXiv}.\footnote{For
traceability, this note corresponds to
\href{https://github.com/LionSR/TNLean/issues/1041}{issue \#1041}. The
underlying question was isolated in
\href{https://github.com/LionSR/TNLean/issues/832}{issue \#832}; the finite
counterexample was added in
\href{https://github.com/LionSR/TNLean/pull/849}{pull request \#849}; and the
positive-semidefinite replacement theorem was added in
\href{https://github.com/LionSR/TNLean/pull/917}{pull request \#917}. The
inverse-map provenance question is recorded in
\href{https://github.com/LionSR/TNLean/issues/4270}{issue \#4270}; the
preceding audit and obstruction verdict are recorded in
\href{https://github.com/LionSR/TNLean/issues/3593}{issue \#3593}. The source
passage is \path{Papers/1606.00608/MPDO-22-12-17-2.tex:1394--1499}.}

The paper constructs positive semidefinite neighboring operators
$\eta_{k,h}$ and defines the trace matrix $T$ by
\begin{align}
    T_{k,h}
        &= \tr(\eta_{k,h}).
    \label{eq:cpgsv17_trace_matrix_definition}
\end{align}
Thus $T$ is entrywise nonnegative. Under the strong area law, the source argues
that $T$ is primitive, meaning that some power of $T$ has all entries strictly
positive. Under zero correlation length, it obtains
\begin{align}
    \tr(T^N)
        &= \tr(T)
    \label{eq:cpgsv17_constant_power_traces}
\end{align}
for every positive integer $N$, with
\begin{align}
    \tr(T)
        &= 1
    \label{eq:cpgsv17_trace_normalization}
\end{align}
after normalization. The proof then invokes fixed-point theory for primitive
nonnegative matrices and concludes that $T$ has one eigenvalue equal to $1$
and all other eigenvalues equal to $0$. From this spectral conclusion, it
asserts the factorization
\begin{align}
    T_{k,h}
        &= a_k b_h.
    \label{eq:cpgsv17_source_outer_product}
\end{align}

The last implication is the point at issue. Primitivity and
Eq.~\ref{eq:cpgsv17_constant_power_traces} control the nonzero eigenvalues of
$T$, but they do not exclude nilpotent Jordan structure on the generalized
eigenspace for the eigenvalue $0$. Such structure is invisible to traces of
positive powers and can prevent $T$ from having rank one.

\section{The counterexample}
\label{sec:cpgsv17_counterexample}

Consider the real $3\times 3$ matrix
\begin{align}
    T
        &=
        \begin{pmatrix}
            \tfrac12 & \tfrac12 & 0\\
            \tfrac16 & \tfrac16 & \tfrac23\\
            \tfrac13 & \tfrac13 & \tfrac13
        \end{pmatrix}.
    \label{eq:cpgsv17_counterexample_matrix}
\end{align}
All entries are nonnegative, and direct multiplication gives
\begin{align}
    T^2
        &= P,
    \label{eq:cpgsv17_counterexample_square}\\
    P
        &:= \frac13
        \begin{pmatrix}
            1&1&1\\
            1&1&1\\
            1&1&1
        \end{pmatrix}.
    \label{eq:cpgsv17_perron_projector}
\end{align}
Hence $T$ is primitive because $T^2$ has all entries strictly positive.
Moreover,
\begin{align}
    P^2
        &= P,
    \label{eq:cpgsv17_projector_idempotence}\\
    PT
        &= P,
    \label{eq:cpgsv17_projector_left_absorption}\\
    TP
        &= P.
    \label{eq:cpgsv17_projector_right_absorption}
\end{align}
Equations~\ref{eq:cpgsv17_counterexample_square}--%
\ref{eq:cpgsv17_projector_right_absorption} imply
\begin{align}
    T^N
        &= P
        \qquad (N\geq 2).
    \label{eq:cpgsv17_counterexample_power_stabilization}
\end{align}
Since
\begin{align}
    \tr(T)
        &= \tr(P)
         = 1,
    \label{eq:cpgsv17_counterexample_trace_one}
\end{align}
it follows that
\begin{align}
    \tr(T^N)
        &= \tr(T)
         = 1
        \qquad (N\geq 1).
    \label{eq:cpgsv17_counterexample_constant_traces}
\end{align}

Nevertheless, $T$ is not an outer product. If $T=ab^{\top}$, then
$T_{1,1}=1/2$ implies $a_1\neq 0$. The equality $T_{1,3}=0$ then forces
$b_3=0$, contradicting $T_{2,3}=2/3$. Thus the universal matrix implication
\begin{align}
    \left(T\ \text{primitive and }
        \tr(T^N)=\tr(T)\ \text{for every }N\geq 1\right)
        &\Longrightarrow \operatorname{rank}(T)=1
    \label{eq:cpgsv17_false_matrix_implication}
\end{align}
is false for finite entrywise nonnegative real matrices.
\footnote{The formal counterexample is
\path{TNLean/Archive/PerronFrobeniusRankOneCounterexample.lean}, especially
\leanid{counterexample} and
\leanid{counterexample_with_rectangular_pairing}. It is also summarized in
\path{docs/counterexamples.md:19--27}.}

The matrix in Eq.~\ref{eq:cpgsv17_counterexample_matrix} has the decomposition
\begin{align}
    T
        &= P+\frac16xy^{\top},
    \label{eq:cpgsv17_nilpotent_perturbation}\\
    x
        &:= (1,-1,0)^{\top},
    \label{eq:cpgsv17_nilpotent_left_vector}\\
    y
        &:= (1,1,-2)^{\top}.
    \label{eq:cpgsv17_nilpotent_right_vector}
\end{align}
The perturbation satisfies
\begin{align}
    (xy^{\top})^2
        &= 0,
    \label{eq:cpgsv17_perturbation_nilpotent}\\
    Pxy^{\top}
        &= 0,
    \label{eq:cpgsv17_perturbation_left_annihilation}\\
    xy^{\top}P
        &= 0.
    \label{eq:cpgsv17_perturbation_right_annihilation}
\end{align}
It therefore leaves all traces of positive powers unchanged from the Perron
projector while increasing the rank of $T$.

The same matrix has a genuinely rectangular realization of the pairing in
lines 1490--1497 of Appendix C.2 of
Ref.~\cite{Cirac2016MPDO_arXiv}. Define
\begin{align}
    L
        &:=
        \begin{pmatrix}
            \tfrac13&\tfrac13&\tfrac13\\
            \tfrac16&\tfrac16&-\tfrac13
        \end{pmatrix},
    \label{eq:cpgsv17_rectangular_left_matrix}\\
    R
        &:=
        \begin{pmatrix}
            1&1\\
            1&-1\\
            1&0
        \end{pmatrix}.
    \label{eq:cpgsv17_rectangular_right_matrix}
\end{align}
Direct multiplication gives
\begin{align}
    RL
        &= T,
    \label{eq:cpgsv17_rectangular_product_trace_matrix}\\
    LR
        &=
        \begin{pmatrix}
            1&0\\
            0&0
        \end{pmatrix},
    \label{eq:cpgsv17_rectangular_product_projection}\\
    (LR)^2
        &= LR.
    \label{eq:cpgsv17_rectangular_idempotence}
\end{align}
Thus the counterexample satisfies not only the scalar trace identities but
also the full rectangular idempotence underlying the displayed
zero-correlation-length equation at source lines 1490--1493 of
Ref.~\cite{Cirac2016MPDO_arXiv}. Nevertheless, $T$ is not an outer product, so
that source equation does not justify the rank-one assertion at lines
1498--1499.

\section{What the operator-valued ZCL identity implies}
\label{sec:cpgsv17_operator_zcl}

The source uses more than the trace equalities before passing to
Perron--Frobenius theory. With $Le_k=l_k$ and
$(Rv)_k=(r_k\mid v)$, one has
\begin{align}
    T
        &= RL.
    \label{eq:cpgsv17_pairing_trace_product}
\end{align}
The displayed ZCL identity at lines 1489--1493 of
Ref.~\cite{Cirac2016MPDO_arXiv} is
\begin{align}
    LR
        &= L(RL)R
         = (LR)^2.
    \label{eq:cpgsv17_source_pairing_idempotence}
\end{align}
This identity makes $LR$ idempotent. Multiplication on the left by $R$ and on
the right by $L$ gives only
\begin{align}
    R(LR)L
        &= R(LR)^2L,
    \label{eq:cpgsv17_rectangular_multiplication_step}\\
    T^2
        &= T^3,
    \label{eq:cpgsv17_square_equals_cube}
\end{align}
not $T=T^2$. The rectangular matrices in
Eqs.~\ref{eq:cpgsv17_rectangular_left_matrix} and
\ref{eq:cpgsv17_rectangular_right_matrix} satisfy the full hypotheses
$T=RL$ and $(LR)^2=LR$, while
\begin{align}
    T^2
        &= T^3
         = P.
    \label{eq:cpgsv17_counterexample_square_cube}
\end{align}
Thus the operator identity alone does not remove the nilpotent part of $T$;
one needs a further property of the inverse-map families $l_k$ and $r_k$.

There is nevertheless a direct rectangular proof of the trace-power statement
at source lines 1494--1497 of
Ref.~\cite{Cirac2016MPDO_arXiv}. Cyclicity of trace gives
\begin{align}
    \tr(T)
        &= \tr(RL)
    \label{eq:cpgsv17_trace_rectangular_product}\\
        &= \tr(LR)
    \label{eq:cpgsv17_trace_cyclic_swap}\\
        &= \tr((LR)^2)
    \label{eq:cpgsv17_trace_pairing_square}\\
        &= \tr(T^2).
    \label{eq:cpgsv17_trace_equals_trace_square}
\end{align}
Together with Eq.~\ref{eq:cpgsv17_square_equals_cube}, this implies
\begin{align}
    T^N
        &= T^2
        \qquad (N\geq 2),
    \label{eq:cpgsv17_power_stabilization_from_square_cube}
\end{align}
and hence
\begin{align}
    \tr(T^N)
        &= \tr(T)
        \qquad (N\geq 1).
    \label{eq:cpgsv17_rectangular_constant_traces}
\end{align}
\footnote{\sloppy Formal declaration:
{\scriptsize\leanid{Matrix.trace_pow_eq_trace_of_rectangular_idempotent}} in
\doclink{TNLean/Algebra/PerronFrobenius/RankOne}.}

The exact sufficient conclusion is transparent. If the inverse-map
construction gave
\begin{align}
    T^2
        &= T,
    \label{eq:cpgsv17_required_trace_matrix_idempotence}
\end{align}
then $T$ itself would be idempotent. An idempotent real matrix has trace equal
to the dimension of its range. Therefore trace one makes the range
one-dimensional and gives $T=ab^{\top}$. This argument is formalized by
\leanid{Matrix.hasRankOneFactorization_of_mul_self_eq_self}. The remaining
source-faithful task is to derive
Eq.~\ref{eq:cpgsv17_required_trace_matrix_idempotence}, or another condition
excluding the generalized zero-eigenspace, from a property of the concrete
inverse-map tensors not currently recorded in
\leanid{MPOTensor.ExplicitEtaOperators}.

\subsection{The pairing-idempotence route}
\label{sec:cpgsv17_pairing_idempotence_route}

The formal development derives $T^2=T$ from
Eq.~\ref{eq:cpgsv17_source_pairing_idempotence} under one additional
hypothesis on the inverse-map tensors. Suppose the family $(l_k)_k$ is linearly
independent, so that $L$ is injective. Then
\begin{align}
    LT^2
        &= (LR)^2L
    \label{eq:cpgsv17_injective_pairing_first_step}\\
        &= (LR)L
    \label{eq:cpgsv17_injective_pairing_second_step}\\
        &= LT.
    \label{eq:cpgsv17_injective_pairing_cancellation_input}
\end{align}
Cancelling $L$ on the left yields
\begin{align}
    T^2
        &= T.
    \label{eq:cpgsv17_injective_pairing_idempotence}
\end{align}
If instead the functionals $(r_k)_k$ are linearly independent, the same
argument in the dual space gives the idempotence of $T^{\top}$ and hence of
$T$. Together with $\tr(T)=1$ and the idempotent trace-one criterion, this
recovers the outer-product factorization asserted by Lemma~C.5 of
Ref.~\cite{Cirac2016MPDO_arXiv}.\footnote{\sloppy The formal statements are
{\scriptsize\leanid{Matrix.mul_self_eq_self_of_pairing_idempotent}},
{\scriptsize\leanid{Matrix.mul_self_eq_self_of_pairing_idempotent_dual}}, and
{\scriptsize\leanid{Matrix.hasRankOneFactorization_of_pairing_idempotent}} in
\doclink{TNLean/Algebra/PerronFrobenius/RankOne}, added in \ghpr{3685}.}

The linear-independence hypothesis is a scope restriction relative to the
source. Lemma~C.5 neither assumes nor derives it. The construction preceding
the lemma uses injectivity of $\mathcal K$ to prove that the sector
decomposition of the trace matrix is trivial, meaning that $T$ is primitive
\cite[Appendix~C.2, lines~1457--1470]{Cirac2016MPDO_arXiv}. Primitivity
constrains only the support pattern of the nonnegative matrix $T$. A primitive
$T$ can arise from linearly dependent $l_k$, so the cited passage does not
make either family $(l_k)_k$ or $(r_k)_k$ linearly independent. The proof of
Lemma~C.5 then uses primitivity in the Perron--Frobenius trace-power step, not
independence of these families. The formal lemmas therefore state independence
as an explicit hypothesis. Whether the inverse-map construction supplies this
property, for instance through injectivity of $\mathcal K$, remains to be
established at the matrix product density operator call site.

\subsection{The sector pairing at the matrix product density operator side}
\label{sec:cpgsv17_sector_pairing_mpdo}

The formal development records the paper's pairing on the matrix product
density operator side. In an abstract real vector space, let $|l_k)$ be the
closed sector tensors and $(r_k|$ the corresponding functionals. It assumes
the two displayed identities at lines 1478--1493 of Appendix C.2 of
Ref.~\cite{Cirac2016MPDO_arXiv}:
\begin{align}
    T_{k,h}
        &= (r_k\mid l_h),
    \label{eq:cpgsv17_sector_pairing_entries}\\
    \sum_k |l_k)(r_k|
        &= \sum_{k,h}T_{k,h}|l_k)(r_h|.
    \label{eq:cpgsv17_sector_pairing_operator_identity}
\end{align}
Define the pairing operator by
\begin{align}
    M
        &:= \sum_k |l_k)(r_k|.
    \label{eq:cpgsv17_pairing_operator_definition}
\end{align}
Equations~\ref{eq:cpgsv17_sector_pairing_entries} and
\ref{eq:cpgsv17_sector_pairing_operator_identity} make $M$ idempotent without
requiring independence of $(l_k)_k$ or $(r_k)_k$. Pairing this idempotence with
$R$ on the left and $L$ on the right, as in
Section~\ref{sec:cpgsv17_operator_zcl}, gives $T^2=T^3$ unconditionally.
Restricting $M$ to the finite-dimensional span of $(l_k)_k$ identifies
$\tr(T)$ with $\tr(T^2)$, so every positive power of $T$ has the same trace as
$T$. This recovers the trace identity at lines 1494--1497 of Appendix C.2 of
Ref.~\cite{Cirac2016MPDO_arXiv} unconditionally on the matrix product density
operator side.
\footnote{\sloppy The formal declarations are
{\scriptsize\leanid{MPOTensor.SectorPairingData.pairing_sq_eq_pairing_cube}}
and
{\scriptsize\leanid{MPOTensor.SectorPairingData.tracePowersConstant}} in
\doclink{TNLean/MPS/MPDO/SimpleLocalStructure}, resting on the general matrix
statements
{\scriptsize\leanid{Matrix.pow_two_eq_pow_three_of_pairing_idempotent}},
{\scriptsize\leanid{Matrix.trace_eq_trace_sq_of_pairing_idempotent}},
{\scriptsize\leanid{Matrix.pow_eq_pow_two_of_pow_two_eq_pow_three}}, and
{\scriptsize\leanid{Matrix.tracePowersConstant_of_pairing_idempotent}} in
\doclink{TNLean/Algebra/PerronFrobenius/RankOne}, added for \ghissue{3714}.}

If the vectors $(l_k)_k$ are also linearly independent, then the coefficient
map $L$ is injective and
\begin{align}
    LT^2
        &= M^2L
    \label{eq:cpgsv17_sector_pairing_injective_first_step}\\
        &= ML
    \label{eq:cpgsv17_sector_pairing_injective_second_step}\\
        &= LT.
    \label{eq:cpgsv17_sector_pairing_injective_cancellation}
\end{align}
Thus $T^2=T$, which is strictly stronger than $T^2=T^3$. Trace one and
idempotence then give the rank-one factorization directly. Primitivity
separately makes each $|l_k)$ nonzero, so an independent family of subspaces
carrying these vectors suffices to derive their linear independence.
\footnote{\sloppy The formal declarations are
{\scriptsize\leanid{MPOTensor.SectorPairingData.pairing_operator_idempotent}},
{\scriptsize\leanid{MPOTensor.SectorPairingData.mul_self_eq_self}},
{\scriptsize\leanid{MPOTensor.SectorPairingData.l_ne_zero}},
{\scriptsize\leanid{MPOTensor.SectorPairingData.linearIndependent_l}},
{\scriptsize\leanid{MPOTensor.sal_zcl_implies_rank_one_T_of_pairing_idempotent}},
and
{\scriptsize\leanid{MPOTensor.sal_zcl_implies_rank_one_T_of_sector_supports}}
in \doclink{TNLean/MPS/MPDO/SimpleLocalStructure}.}

The inverse-map construction has been connected to this argument without
additional hypotheses. The closed tensors $(l_k)_k$ and $(r_k)_k$ are
constructed from an injective simple tensor, the complex matrix is defined by
\begin{align}
    T_{k,h}
        &= (r_k\mid l_h)
         = \tr(\eta_{k,h}),
    \label{eq:cpgsv17_closed_sector_trace_matrix}
\end{align}
and its pairing operator is identified with the physical-trace transfer.
Source zero correlation length therefore supplies a number $\lambda>0$ for
which the normalized pairing operator is idempotent. The rectangular
factorizations
\begin{align}
    S
        &= LR,
    \label{eq:cpgsv17_inverse_map_physical_transfer}\\
    T
        &= RL
    \label{eq:cpgsv17_inverse_map_trace_matrix}
\end{align}
then prove
\begin{align}
    (\lambda^{-1}T)^2
        &= (\lambda^{-1}T)^3,
    \label{eq:cpgsv17_normalized_square_cube}\\
    \tr((\lambda^{-1}T)^N)
        &= \tr(\lambda^{-1}T)
        \qquad (N\geq 1).
    \label{eq:cpgsv17_normalized_constant_traces}
\end{align}
\footnote{\sloppy Formal declarations:
{\scriptsize\leanid{MPOTensor.closedSectorTraceMatrix_apply}} and
{\scriptsize\leanid{MPOTensor.closedSectorTraceMatrix_normalized_relations}}
in \doclink{TNLean/MPS/MPDO/SectorPairingTransfer}.}

The normalized identities can also be placed on the same real trace matrix
for which the strong-area-law argument proves primitivity. After discarding
only zero-weight indices and choosing coherent sector phases, define
\begin{align}
    T^*_{k,h}
        &:= \operatorname{Re}\tr(\eta_{k,h})
        \qquad (p_kp_h\neq 0).
    \label{eq:cpgsv17_active_real_trace_matrix}
\end{align}
The left tensor of every zero-weight sector vanishes. Consequently, the
physical-trace transfer retains the active rectangular factorization $LQ$,
while positivity of the neighboring operators identifies $QL$ with the
complexification of $T^*$. The same scalar $\lambda>0$ therefore gives
\begin{align}
    T^*
        &\ \text{primitive},
    \label{eq:cpgsv17_active_trace_matrix_primitive}\\
    (\lambda^{-1}T^*)^2
        &= (\lambda^{-1}T^*)^3,
    \label{eq:cpgsv17_active_trace_matrix_square_cube}\\
    \tr((\lambda^{-1}T^*)^N)
        &= \tr(\lambda^{-1}T^*)
        \qquad (N\geq 1).
    \label{eq:cpgsv17_active_trace_matrix_constant_traces}
\end{align}
\footnote{\sloppy Formal declarations:
{\scriptsize\leanid{MPOTensor.PhysicalSectorFactorization.activeSectorTraceMatrix_normalized_relations_of_isSourceZCL}},
{\scriptsize\leanid{MPOTensor.exists_rephased_inverseMap_activeSectorTraceMatrix_normalized_relations}},
and
{\scriptsize\leanid{MPOTensor.exists_activeTraceMatrix_relations_of_isSAL_of_isSourceZCL}}
in \doclink{TNLean/MPS/MPDO/InverseMapActiveSectorZCL}.}

Write $\mathcal B_{\mathcal K}$ for the physical-trace transfer of
$\mathcal K$. If SAL is accompanied by the paper's literal identity
\begin{align}
    \mathcal B_{\mathcal K}^2
        &= \mathcal B_{\mathcal K},
    \label{eq:cpgsv17_literal_zcl_identity}
\end{align}
then the one-site normalization clause in SAL first gives
$\mathcal B_{\mathcal K}\neq 0$. Thus literal ZCL implies source ZCL, and the
same primitive active trace matrix and normalized relations follow without
additional hypotheses. This specialization is formalized by
\leanid{MPOTensor.exists_activeTraceMatrix_relations_of_isSAL_of_literal_ZCL}
in \doclink{TNLean/MPS/MPDO/InverseMapActiveSectorZCL}.

This conclusion is the scale-invariant consequence of the displayed operator
identity. It does not imply that the normalized trace matrix is idempotent or
semisimple, nor that it has rank one. The remaining pairing route must prove
linear independence of one of the two closed tensor families or supply another
argument excluding the generalized zero-eigenspace. The direct sum in the
source's sector splitting at lines 1436--1448 of
Ref.~\cite{Cirac2016MPDO_arXiv} concerns physical indices. After those indices
are contracted, the closed tensors live in a common bond space. The displayed
equations alone therefore do not place them in corresponding members of an
independent family of subspaces or otherwise prove their independence.

\subsection{Virtual spanning does not eliminate the nilpotent part}
\label{sec:cpgsv17_virtual_spanning}

Injectivity gives more than primitivity in the present formal development: the
sector virtual matrices span the full virtual matrix algebra. This additional
conclusion still does not imply idempotence of the opposite product. Consider
\begin{align}
    L
        &:=
        \begin{pmatrix}
            \tfrac16&\tfrac16&\tfrac13&\tfrac13\\
            0&\tfrac16&\tfrac16&-\tfrac13
        \end{pmatrix},
    \label{eq:cpgsv17_spanning_left_matrix}\\
    Q
        &:=
        \begin{pmatrix}
            1&1\\
            1&1\\
            1&-1\\
            1&0
        \end{pmatrix}.
    \label{eq:cpgsv17_spanning_right_matrix}
\end{align}
Write $l_k$ for the $k$th column of $L$ and $r_k$ for the $k$th row of $Q$.
Direct multiplication gives
\begin{align}
    LQ
        &=
        \begin{pmatrix}
            1&0\\
            0&0
        \end{pmatrix},
    \label{eq:cpgsv17_spanning_physical_projection}\\
    T=QL
        &=
        \begin{pmatrix}
            \tfrac16&\tfrac13&\tfrac12&0\\
            \tfrac16&\tfrac13&\tfrac12&0\\
            \tfrac16&0&\tfrac16&\tfrac23\\
            \tfrac16&\tfrac16&\tfrac13&\tfrac13
        \end{pmatrix}.
    \label{eq:cpgsv17_spanning_trace_matrix}
\end{align}
The square of $T$ is
\begin{align}
    T^2
        &=
        \begin{pmatrix}
            \tfrac16&\tfrac16&\tfrac13&\tfrac13\\
            \tfrac16&\tfrac16&\tfrac13&\tfrac13\\
            \tfrac16&\tfrac16&\tfrac13&\tfrac13\\
            \tfrac16&\tfrac16&\tfrac13&\tfrac13
        \end{pmatrix}.
    \label{eq:cpgsv17_spanning_trace_matrix_square}
\end{align}
Thus $T$ is primitive, $\tr(T)=1$, $T^2=T^3$, and $T\neq T^2$. Equivalently,
the missing expression is nonzero:
\begin{align}
    Q(\Id-LQ)L
        &= T-T^2
         \neq 0.
    \label{eq:cpgsv17_spanning_nonzero_remainder}
\end{align}
Moreover, the four matrices $l_kr_k$ span $M_2(\mathbb R)$. After vectorizing
them in the order $(00,01,10,11)$, the determinant of the resulting
four-by-four matrix is
\begin{align}
    \det V
        &= \frac{1}{108}.
    \label{eq:cpgsv17_spanning_vectorization_determinant}
\end{align}
Consequently, full virtual spanning does not exclude the generalized
zero-eigenspace.

This example is realized at the level of physical slices by the diagonal
matrix product operator tensor
\begin{align}
    K^{ij}
        &= \delta_{ij}\,l_i r_i.
    \label{eq:cpgsv17_spanning_tensor_slices}
\end{align}
Its doubled-index tensor is injective because the four nonzero slices span
$M_2(\mathbb C)$, and its physical-trace transfer is
\begin{align}
    \sum_i K^{ii}
        &= LQ.
    \label{eq:cpgsv17_spanning_physical_trace_transfer}
\end{align}
It therefore satisfies source zero correlation length. These assertions,
together with primitivity, trace normalization, the spanning identity, and
the nonzero nilpotent remainder, are proved in
\path{TNLean/Archive/PerronFrobeniusVirtualSpanningCounterexample.lean}.
\footnote{\sloppy Combined formal declaration:
{\scriptsize\leanid{TNLean.Archive.PerronFrobeniusVirtualSpanningCounterexample.injective_sourceZCL_tensor_with_nilpotent_sector_pairing}}.}

The construction does not identify $L$ and $Q$ with the particular closed
tensors selected from a Hayashi decomposition by
\leanid{MPOTensor.zeroWeightReparameterizedInverseMapPhysicalSectorFactorization}.
This distinction is essential. The theorem supplies a directly chosen
rectangular realization; it does not supply the Hayashi structure, normalized
tail entry, inverse-map indices, and rephasing data from which the source
factorization is constructed. It is therefore not a counterexample to the
full statement of Lemma~C.5 of Ref.~\cite{Cirac2016MPDO_arXiv}.

The obstruction instead proves that every presently exported algebraic
consequence---source ZCL, injectivity, primitivity, trace normalization,
rectangular pairing idempotence, constant positive-power traces, and full
virtual spanning---is insufficient. A source-faithful completion must derive
an additional identity from SAL and the provenance of the inverse-map tensors.
In the notation above, the required assertion is
\begin{align}
    Q(\Id-LQ)L
        &= 0,
    \label{eq:cpgsv17_required_rectangular_vanishing}
\end{align}
or any equivalent statement excluding the nilpotent generalized
zero-eigenspace of $QL$.

\subsection{Positive semidefinite neighboring operators do not force the rectangular vanishing}
\label{sec:cpgsv17_psd_neighboring_obstruction}

The preceding obstruction can also be realized inside the physical-sector
factorization with every neighboring operator positive semidefinite. A second,
independent witness has bond dimension two and four one-dimensional physical
sectors. Define
\begin{align}
    L
        &:=
        \begin{pmatrix}
            \frac14&\frac14&\frac14&\frac14\\
            1&-1&1&-1
        \end{pmatrix},
    \label{eq:cpgsv17_psd_witness_left_matrix}\\
    Q
        &:=
        \begin{pmatrix}
            1&\frac18\\
            1&\frac18\\
            1&-\frac18\\
            1&-\frac18
        \end{pmatrix}.
    \label{eq:cpgsv17_psd_witness_right_matrix}
\end{align}
For each sector $k$, take the virtual matrix to be the outer product of the
$k$th column of $L$ with the $k$th row of $Q$. These four matrices span
$\MN{2}$: in the coordinate order $(00,01,10,11)$, the matrix having these
four outer products as columns has determinant
\begin{align}
    \det W
        &= -\frac{1}{64}.
    \label{eq:cpgsv17_psd_witness_spanning_determinant}
\end{align}
Placing them on the four diagonal physical entries therefore gives an
injective tensor with a scalar physical-sector factorization.

The two rectangular products are
\begin{align}
    LQ
        &=
        \begin{pmatrix}
            1&0\\
            0&0
        \end{pmatrix},
    \label{eq:cpgsv17_psd_witness_physical_projection}\\
    QL
        &=
        \begin{pmatrix}
            \frac38&\frac18&\frac38&\frac18\\
            \frac38&\frac18&\frac38&\frac18\\
            \frac18&\frac38&\frac18&\frac38\\
            \frac18&\frac38&\frac18&\frac38
        \end{pmatrix}.
    \label{eq:cpgsv17_psd_witness_trace_matrix}
\end{align}
Hence $LQ$ is a nonzero idempotent physical-trace transfer, while $QL$ is
strictly positive, primitive, and has trace one. Its entries are the scalar
neighboring operators and are therefore positive semidefinite. Nevertheless,
\begin{align}
    (QL)^2
        &\neq QL,
    \label{eq:cpgsv17_psd_witness_nonidempotence}\\
    ((QL)^2)_{0,1}
        &= \frac14
         \neq \frac18
         = (QL)_{0,1}.
    \label{eq:cpgsv17_psd_witness_changed_entry}
\end{align}
Equivalently,
\begin{align}
    Q(\Id-LQ)L
        &= QL-(QL)^2
         \neq 0.
    \label{eq:cpgsv17_psd_witness_nonzero_remainder}
\end{align}
Thus injectivity, full virtual-matrix spanning, positive semidefinite
neighboring operators, primitivity, trace normalization, and the source
zero-correlation-length rectangular identity can all hold without the missing
vanishing.
\footnote{The formal tensor and combined counterexample theorem are in
\doclink{TNLean/MPS/MPDO/ActiveSectorSpanningCounterexample}; see
{\scriptsize\leanid{MPOTensor.ActiveSectorSpanningCounterexample.virtual_spanning_does_not_force_rectangular_remainder_zero}}.}

Let $T:=QL$. The induced finite-chain state is the classical cyclic law
\begin{align}
    p_N(i_1,\ldots,i_N)
        &:= \prod_{n=1}^N T_{i_n,i_{n+1}},
    \label{eq:cpgsv17_classical_cyclic_law}\\
    i_{N+1}
        &:= i_1.
    \label{eq:cpgsv17_cyclic_boundary_index}
\end{align}
The matrix $T$ is doubly stochastic, and $T^2$ has every entry equal to
$1/4$. Let $h$ be the common Shannon entropy, in bits, of a row of $T$; the
four rows are permutations of one another. For every proper nonempty block of
length $L$ in a chain relevant to SAL,
\begin{align}
    H_L
        &= 2+(L-1)h,
    \label{eq:cpgsv17_block_entropy}\\
    H_N
        &= Nh.
    \label{eq:cpgsv17_chain_entropy}
\end{align}
Consequently,
\begin{align}
    I_L
        &= H_L+H_{N-L}-H_N
    \label{eq:cpgsv17_mutual_information_definition}\\
        &= 4-2h,
    \label{eq:cpgsv17_mutual_information_constant}
\end{align}
independently of $L$. Formally, the same conclusion follows without
logarithmic calculations. After one site is traced out, the diagonal state is
an ordinary Markov path, and conditioning on its middle label gives a
four-sector Hayashi decomposition at every admissible cut. The theorem
\leanid{MPOTensor.ActiveSectorSpanningCounterexample.tensor_isSAL} therefore
proves the strong area law for this tensor.

The inverse-map provenance is explicit. The refined Hayashi decomposition has
four one-dimensional sectors of weight $1/4$, with diagonal left and right
density matrices determined by $QL$. Its normalized four-site tail is $LQ$.
The four nonzero physical slices form a basis of the virtual matrix algebra;
computing the dual basis shows that the source inverse map recovers the
displayed closed tensors $L$ and $Q$ exactly. Consequently, its neighboring
operators agree with those of
\leanid{MPOTensor.ActiveSectorSpanningCounterexample.factorization}, and the
same source-selected construction satisfies
\begin{align}
    Q(\Id-LQ)L
        &\neq 0.
    \label{eq:cpgsv17_inverse_map_nonzero_remainder}
\end{align}
These statements are proved in
\leanid{MPOTensor.ActiveSectorSpanningCounterexample.%
inverseMap_provenance_preserves_rectangular_remainder}. Together with SAL,
injectivity, and source ZCL, these facts are assembled in
\leanid{MPOTensor.ActiveSectorSpanningCounterexample.%
tensor_has_sal_sourceZCL_and_non_rankOne_activeTraceMatrix}. The active trace
matrix has no factorization $T_{k,h}=a_kb_h$, while
Eq.~\ref{eq:cpgsv17_inverse_map_nonzero_remainder} holds. Hence the example
proves the full raw-tensor conjunction, including inverse-map provenance.

It does not refute Lemma~C.5 under its complete standing hypotheses
\cite[Appendix~C.2]{Cirac2016MPDO_arXiv}. The lemma lies inside Case~I at
source lines 1374--1381, where $\mathcal K$ is assumed to be an injective
normal tensor. The raw witness is not transfer-normalized, while its normal
representative loses the literal ZCL identity. This normalization
qualification corrects the earlier verdict tracked in \ghissue{4270}.

\section{Renormalization maps and the global normalization boundary}
\label{sec:cpgsv17_global_normalization}

The standing hypotheses of Theorem~4.9 of
Ref.~\cite{Cirac2016MPDO_arXiv} require a simple tensor in biCF with a basis of
normal tensors. They are imposed at source lines 849--893 and restated at
line 1628. They also inherit the global canonical-form convention at source
line 246: every BNT coefficient has modulus at most one, and at least one has
modulus one. The distinction between this single global unit witness and a
stronger sectorwise convention is recorded in
\path{docs/paper-gaps/cpsv16_global_vs_persector_unit_witness.tex}.

Write $\mathcal K$ for the four-sector witness and regard its doubled physical
index as an MPS index. Its four nonzero physical matrices span $\MN{2}$, so
the doubled tensor is injective at one site. Its raw transfer map is
\begin{align}
    \mathcal E_{\mathcal K}(X)
        &=
        \left(2X_{00}+\frac1{32}X_{11}\right)
        \begin{pmatrix}
            \frac18&0\\
            0&2
        \end{pmatrix},
    \label{eq:cpgsv17_raw_transfer_map}\\
    \mathcal E_{\mathcal K}^{2}
        &= \frac5{16}\mathcal E_{\mathcal K}.
    \label{eq:cpgsv17_raw_transfer_quasi_idempotence}
\end{align}
Thus the raw doubled tensor has spectral radius $5/16$. Define
\begin{align}
    A
        &:= \frac4{\sqrt5}\mathcal K,
    \label{eq:cpgsv17_normalized_representative}\\
    \mu
        &:= \frac{\sqrt5}{4}.
    \label{eq:cpgsv17_raw_bnt_coefficient}
\end{align}
Then
\begin{align}
    \mathcal K
        &= \mu A,
    \label{eq:cpgsv17_raw_normalized_decomposition}
\end{align}
and $\mathcal E_A$ has spectrum $\{1,0,0,0\}$ and is idempotent. The singleton
family $\{A\}$ is the normal basis for $\mathcal K$, but its only coefficient
is $\mu<1$. Consequently, the raw tensor $\mathcal K$ does not meet the
source's global unit-weight convention. The normalized representative $A$
does: it is injective and hence normal; its singleton BNT has coefficient one;
and the same one-site spanning identity gives biCF with blocking length one.
Moreover,
\begin{align}
    \mathcal B_A
        &= \frac4{\sqrt5}LQ
    \label{eq:cpgsv17_normalized_physical_transfer_scaled}\\
        &= \frac4{\sqrt5}
        \begin{pmatrix}
            1&0\\
            0&0
        \end{pmatrix}.
    \label{eq:cpgsv17_normalized_physical_transfer_matrix}
\end{align}
This operator is not nilpotent, so the sole BNT element is simple in the sense
used in Theorem~4.9 of Ref.~\cite{Cirac2016MPDO_arXiv}.

These two representatives lie on opposite sides of the source's two literal
normalizations. The raw tensor has
\begin{align}
    \mathcal B_{\mathcal K}
        &= LQ
         = P,
    \label{eq:cpgsv17_raw_physical_transfer_projection}
\end{align}
so it satisfies
\begin{align}
    \mathcal B_{\mathcal K}^2
        &= \mathcal B_{\mathcal K},
    \label{eq:cpgsv17_raw_literal_zcl}
\end{align}
but it lacks a global unit-modulus BNT coefficient. The representative $A$
has the required unit BNT coefficient, but
\begin{align}
    \mathcal B_A^2
        &= \frac4{\sqrt5}\mathcal B_A
    \label{eq:cpgsv17_normalized_scaled_zcl}\\
        &\neq \mathcal B_A.
    \label{eq:cpgsv17_normalized_literal_zcl_failure}
\end{align}
Thus $A$ fails the paper's literal zero-correlation-length diagram. The scalar
absorption at source line 1660 of Ref.~\cite{Cirac2016MPDO_arXiv} occurs only
after that literal ZCL identity has been assumed; it does not transport
literal ZCL or the renormalization maps across scalar rescaling.

The raw tensor nevertheless gives a positive calculation for the
trace-preserving completely positive map condition. Let $\tau:=QL$ and define
\begin{align}
    f(i,j)
        &:= (i\bmod 2)+2\left\lfloor\frac j2\right\rfloor.
    \label{eq:cpgsv17_label_coarsening_function}
\end{align}
The sector matrices $S_i$ obey
\begin{align}
    S_iS_j
        &= \tau_{ij}S_{f(i,j)},
    \label{eq:cpgsv17_sector_matrix_multiplication}\\
    \sum_{(i,j):f(i,j)=k}\tau_{ij}
        &= 1
        \qquad (k\in\{0,1,2,3\}).
    \label{eq:cpgsv17_coarsening_weight_normalization}
\end{align}
Define the refinement and coarse-graining Kraus operators by
\begin{align}
    R_{ij}
        &:= \sqrt{\tau_{ij}}\ket{i,j}\bra{f(i,j)},
    \label{eq:cpgsv17_refinement_kraus_operator}\\
    C_{ij}
        &:= \ket{f(i,j)}\bra{i,j}.
    \label{eq:cpgsv17_coarse_graining_kraus_operator}
\end{align}
These operators resolve their respective input identities. Their maps are
trace-preserving and completely positive, and direct substitution gives
\begin{align}
    \mathsf C[\mathcal K_2(X)]
        &= \mathcal K_1(X),
    \label{eq:cpgsv17_coarse_graining_closure}\\
    \mathsf R[\mathcal K_1(X)]
        &= \mathcal K_2(X)
    \label{eq:cpgsv17_refinement_closure}
\end{align}
for every virtual matrix $X$. Applying each map twice and passing to canonical
blocked coordinates gives the same identities for $\mathcal K^{[2]}$. These
statements are proved by
\leanid{MPOTensor.ActiveSectorSpanningCounterexample.tensor_isRFPViaTS},
\leanid{MPOTensor.ActiveSectorSpanningCounterexample.%
blockTwo_tensor_isRFPViaTS}, and the general blocking theorem
\leanid{MPOTensor.IsRFPViaTS.blockTwo}.

\subsection{Pair-preserving two-to-four-site maps}
\label{sec:cpgsv17_pair_preserving_maps}

There is a sharp trace obstruction to applying the twice-iterated maps while
retaining every fixed pair of outer sector labels. For an arbitrary
physical-sector factorization, define
\begin{align}
    C_{k,h}
        &:= \tr(\eta_{k,h}),
    \label{eq:cpgsv17_neighboring_trace_matrix}\\
    \Theta_{k,h}
        &:=
        \bigoplus_{l,m}
        \eta_{k,l}\otimes\eta_{l,m}\otimes\eta_{m,h}.
    \label{eq:cpgsv17_four_site_neighboring_operator}
\end{align}
Additivity of trace on direct sums and multiplicativity on tensor products
give
\begin{align}
    \tr(\Theta_{k,h})
        &= \sum_{l,m}
        \tr(\eta_{k,l})
        \tr(\eta_{l,m})
        \tr(\eta_{m,h})
    \label{eq:cpgsv17_four_site_trace_expansion}\\
        &= \sum_{l,m}C_{k,l}C_{l,m}C_{m,h}
    \label{eq:cpgsv17_four_site_trace_matrix_sum}\\
        &= (C^3)_{k,h}.
    \label{eq:cpgsv17_four_site_trace_cube}
\end{align}
It follows that a trace-preserving family that sends every $\eta_{k,h}$ to
the corresponding $\Theta_{k,h}$, or every $\Theta_{k,h}$ to the corresponding
$\eta_{k,h}$, forces
\begin{align}
    C
        &= C^3.
    \label{eq:cpgsv17_pair_preserving_trace_condition}
\end{align}
This project-derived result is a \emph{conditional helper}: it identifies the
necessary trace equality for a pair-preserving construction, but it neither
asserts that the source maps preserve such pairs nor constructs a label-mixing
alternative. The trace identity and its two directions are proved in
\leanid{MPOTensor.PhysicalSectorFactorization.%
trace_fourSiteNeighboringOperator},
\leanid{MPOTensor.PhysicalSectorFactorization.%
neighboringTraceMatrix_eq_pow_three_of_pairwise_coarsening}, and
\leanid{MPOTensor.PhysicalSectorFactorization.%
neighboringTraceMatrix_eq_pow_three_of_pairwise_refinement}.

For the selected four-sector factorization of the active-spanning witness,
$C=QL$. Rectangular idempotence of $LQ$ gives
\begin{align}
    C^2
        &= C^3,
    \label{eq:cpgsv17_selected_trace_square_cube}
\end{align}
but the explicit entries satisfy
\begin{align}
    C_{0,1}
        &= \frac18,
    \label{eq:cpgsv17_selected_trace_entry}\\
    (C^2)_{0,1}
        &= \frac14.
    \label{eq:cpgsv17_selected_trace_square_entry}
\end{align}
Hence
\begin{align}
    C^2
        &= C^3,
    \label{eq:cpgsv17_pair_preserving_square_cube}\\
    C
        &\neq C^3.
    \label{eq:cpgsv17_pair_preserving_cube_failure}
\end{align}
Neither direction of a pair-preserving trace-preserving map can therefore
exist for these four selected sectors. This is proved in
\leanid{MPOTensor.ActiveSectorSpanningCounterexample.%
not_exists_pairwise_fourSite_coarsening} and
\leanid{MPOTensor.ActiveSectorSpanningCounterexample.%
not_exists_pairwise_fourSite_refinement}.

This exclusion is \emph{proof-path drift}, not an obstruction to
renormalization. The trace-preserving completely positive maps above use the
many-to-one rule $f(i,j)$ and therefore coarsen the selected labels; after
iteration, they still prove the renormalization equations for
$\mathcal K^{[2]}$. They are not among the fixed-$(k,h)$ families excluded by
the trace calculation. The one-sector factorization likewise has no such
four-sector obstruction. Thus the calculation rules out only the selected
pair-preserving route. The explicit label-coarsening maps belong to this
example; they are not the general construction used in the source proof.
Proposition~\texttt{3to5} of Ref.~\cite{Cirac2016MPDO_arXiv} instead
constructs its maps from normalized rank-one neighboring traces, and the proof
of Proposition~\texttt{prop2to5} only adds the projection onto the outer BNT
subspaces before applying those maps.

This positive result is not invariant under the normalization above. Define
\begin{align}
    P
        &:= \operatorname{diag}(1,0).
    \label{eq:cpgsv17_normalized_coordinate_projection}
\end{align}
The physical-trace transfers of the normalized representative and its
two-site blocking are
\begin{align}
    \mathcal B_A
        &= \frac4{\sqrt5}P,
    \label{eq:cpgsv17_normalized_transfer_projection}\\
    \mathcal B_{A^{[2]}}
        &= \mathcal B_A^2
    \label{eq:cpgsv17_blocked_transfer_square}\\
        &= \frac{16}{5}P.
    \label{eq:cpgsv17_blocked_transfer_projection}
\end{align}
Definition~4.1 of Ref.~\cite{Cirac2016MPDO_arXiv} would force
$\mathcal B_{A^{[2]}}$ to be literally idempotent because trace preservation
identifies the traces of its one-site and two-site closures for every virtual
boundary matrix. However,
\begin{align}
    \left(\frac{16}{5}P\right)^2
        &\neq \frac{16}{5}P.
    \label{eq:cpgsv17_blocked_transfer_nonidempotence}
\end{align}
Hence $A^{[2]}$ is not a renormalization fixed point. This is proved by
\leanid{MPOTensor.ActiveSectorSpanningCounterexample.%
blockTwo_normalizedTensor_not_isRFPViaTS}.
\footnote{The source-hypothesis and normalization-sensitive obstruction are
tracked in \ghissue{5388}.}

There are two distinct zero-correlation-length statements. TNLean's
\leanid{MPOTensor.IsSourceZCL} interpretation is scale-invariant: it asks for
\begin{align}
    \mathcal B_A^2
        &= \lambda\mathcal B_A
        \qquad (\lambda>0).
    \label{eq:cpgsv17_scale_invariant_zcl}
\end{align}
The paper's diagrammatic identity is literal and selects a fixed scalar
representative. Both the strong area law and TNLean's up-to-scalar predicate
are unchanged by multiplying the tensor by $4/\sqrt5$, but the paper's diagram
is not. Hence $A$ is a counterexample only to the broadened implication
obtained by replacing the literal diagram with
\leanid{MPOTensor.IsSourceZCL}; it is not an unqualified counterexample to
Theorem~4.9 of Ref.~\cite{Cirac2016MPDO_arXiv}. This statement concerns only
the four-sector witness $A$. The scalar repeated-copy counterexample below is
an exact counterexample to the printed implication (iv)$\Rightarrow$(v).

The scalar identities, the raw renormalization maps, and the failure of
Definition~4.1 for $A^{[2]}$ are kernel-checked. This example leaves the source
(ii)$\Rightarrow$(v) boundary unsettled: $\mathcal K$ satisfies literal ZCL
but not the global unit-weight convention, while $A$ satisfies the global
convention but not literal ZCL. Packaging the normality, singleton BNT, biCF,
simplicity, and SAL facts for $A$ remains useful, but cannot turn its
up-to-scalar ZCL relation into the paper's literal diagram. The distinct raw
copy example below settles (iv)$\Rightarrow$(v) without asserting
condition~(ii).

\section{A positive-semidefinite route}
\label{sec:cpgsv17_psd_route}

The formal development proves a corrected sufficient matrix theorem. Let $T$
be a finite real square matrix. If $T$ is positive semidefinite,
$\tr(T)=1$, and
\begin{align}
    \tr(T^N)
        &= \tr(T)
        \qquad (N\geq 1),
    \label{eq:cpgsv17_psd_route_trace_hypothesis}
\end{align}
then $T$ has an outer-product factorization
\begin{align}
    T
        &= ab^{\top}.
    \label{eq:cpgsv17_psd_route_factorization}
\end{align}
\footnote{The formal statements are
\leanid{Matrix.PosSemidef.trace_powers_constant_implies_rank_one} and
\leanid{Matrix.primitive_trace_powers_constant_implies_rank_one_of_posSemidef}
in \path{TNLean/Algebra/PerronFrobenius/RankOne.lean}.}
Primitivity is not needed for this corrected matrix theorem once positive
semidefiniteness and trace normalization are assumed.

The proof is spectral. Positive semidefiniteness makes $T$ Hermitian and
diagonalizable with nonnegative real eigenvalues $\lambda_i$. Trace
normalization gives
\begin{align}
    \sum_i\lambda_i
        &= 1.
    \label{eq:cpgsv17_psd_eigenvalue_sum}
\end{align}
The trace-power hypothesis at $N=2$ gives
\begin{align}
    \sum_i\lambda_i^2
        &= 1.
    \label{eq:cpgsv17_psd_eigenvalue_square_sum}
\end{align}
For nonnegative numbers with sum one,
Eq.~\ref{eq:cpgsv17_psd_eigenvalue_square_sum} forces exactly one $\lambda_i$
to equal $1$ and all others to equal $0$. Therefore the diagonal form has rank
one, and unitary conjugation gives Eq.~\ref{eq:cpgsv17_psd_route_factorization}.

\section{Audit of the positive-semidefinite route}
\label{sec:cpgsv17_psd_audit}

The neighboring-operator construction displayed in the source does not by
itself exhibit $T$ as a Gram matrix. Immediately before Lemma~C.5, the paper
writes
\begin{align}
    T_{k,h}
        &= (r_k\mid l_h),
    \label{eq:cpgsv17_source_cross_pairing}
\end{align}
with two different families $(r_k)_k$ and $(l_h)_h$
\cite[Appendix~C.2, lines~1473--1482]{Cirac2016MPDO_arXiv}. No equality or
adjoint relation between these families is stated. Consequently,
Eq.~\ref{eq:cpgsv17_source_cross_pairing} is a cross-pairing rather than a Gram
representation and gives neither symmetry nor positive semidefiniteness of
$T$.

The present type of explicit neighboring operators records only
\begin{align}
    \eta_{k,h}
        &\geq 0
        \qquad\text{for each ordered pair }(k,h).
    \label{eq:cpgsv17_neighboring_operator_positivity}
\end{align}
This hypothesis implies
\begin{align}
    T_{k,h}
        &\geq 0
        \qquad\text{entrywise},
    \label{eq:cpgsv17_trace_matrix_entrywise_nonnegative}
\end{align}
but contains no relation between $\eta_{k,h}$ and $\eta_{h,k}$. Thus the type
alone does not imply matrix-level positive semidefiniteness.

There is one special positive-semidefinite case in the formal development. The
sector-reduced Hayashi--Markov family is defined by
\begin{align}
    \eta_{k,h}
        &:=
        \tr_C(\rho_{b_2C}^{(k)})
        \otimes
        \tr_A(\rho_{Ab_1}^{(h)}).
    \label{eq:cpgsv17_hayashi_markov_neighboring_operator}
\end{align}
Since each reduced density matrix has trace one, its trace matrix is the
all-ones matrix. The theorem
\leanid{MPOTensor.ExplicitEtaOperators.traceMatrixRe_ofHayashiMarkov_posSemidef}
proves its positive semidefiniteness. This does not discharge Lemma~C.5 of
Ref.~\cite{Cirac2016MPDO_arXiv}: the sector-reduced family has not been
identified with the inverse-map family constructed at source lines
1413--1455, and the trace of its all-ones matrix is the number of sectors, not
one unless there is a single sector.

The positive-semidefinite route must therefore use more of the matrix product
density operator construction than the separate inequalities in
Eq.~\ref{eq:cpgsv17_neighboring_operator_positivity}. One possible completion
is to derive an adjoint or Gram relation from the inverse-map tensors.
Positive semidefiniteness of the source trace matrix is not established by the
cited passage.

This corrected theorem does not refute the intended matrix product density
operator result; it supplies one sufficient matrix-level hypothesis. For this
route, the open question is whether the trace matrix
$T_{k,h}=\tr(\eta_{k,h})$ is Hermitian or positive semidefinite. The
pairing-idempotence route of
Section~\ref{sec:cpgsv17_pairing_idempotence_route} avoids that hypothesis but
instead requires the concrete sector pairing, the zero-correlation-length
operator identity, and linear independence. Positivity of each $\eta_{k,h}$
gives only entrywise nonnegativity of $T$, which implies neither matrix-level
positive semidefiniteness nor Hermitian symmetry.
\footnote{The entrywise nonnegativity result is
{\scriptsize\leanid{MPOTensor.ExplicitEtaOperators.traceMatrixRe_nonneg}}. The
positive-semidefinite reformulation for the rank-one conclusion is
{\scriptsize\leanid{MPOTensor.sal_zcl_implies_rank_one_T_of_posSemidef}}. The
corresponding blueprint entries have labels
\texttt{\detokenize{thm:mpdo\_explicit\_eta\_trace\_matrix\_nonneg}} and
\texttt{\detokenize{thm:psd\_trace\_powers\_rank\_one\_t}} in
\path{blueprint/src/chapter/ch21_mpdo_rfp_simple_local_primitivity.tex}.}

\section{Audit of the inverse-map provenance}
\label{sec:cpgsv17_inverse_map_audit}

The remaining possibility was that the particular inverse-map tensors used in
Appendix C.2 of Ref.~\cite{Cirac2016MPDO_arXiv} satisfy an additional identity
not visible in the rectangular pairing equation. An audit of the complete
source passage does not reveal such an identity.

In the proof of Lemma~C.4 at source lines 1406--1471, the strong area law is
used to obtain the physical-sector factorization, positive neighboring
operators, and primitivity of the trace matrix
\cite[Appendix~C.2]{Cirac2016MPDO_arXiv}. The last conclusion uses injectivity
to rule out more than one primitive block. The argument does not prove an
adjoint relation between the closed left and right tensors, their linear
independence, or the vanishing
\begin{align}
    Q(\Id-LQ)L
        &= 0.
    \label{eq:cpgsv17_audit_required_vanishing}
\end{align}
The physical direct sum separates the uncontracted physical sector spaces.
After the physical legs are closed at source lines 1473--1482, the vectors
$l_h$ and functionals $r_k$ lie in common bond spaces, and the source gives
only $T_{k,h}=(r_k\mid l_h)$.

The proof of Lemma~C.5 at source lines 1489--1499 then uses zero correlation
length to obtain the pairing-operator identity and its trace-power consequence
\cite[Appendix~C.2]{Cirac2016MPDO_arXiv}. No further property of the inverse
map occurs between this identity and the rank-one conclusion. At the
common-matrix level, this leaves the gap described in
Section~\ref{sec:cpgsv17_assertion}. The separate Case~I assumption that
$\mathcal K$ is an injective normal tensor must nevertheless be retained when
judging the full lemma.

The cyclic finite-chain form at source lines 1581--1593 does not repair this
step. It is obtained later from the same positive neighboring family and
states a direct sum of cyclic products
\cite[Appendix~C.2]{Cirac2016MPDO_arXiv}. Its formal, scalar-preserving version
is
\leanid{MPOTensor.EtaLocalStructureData.%
exists_positive_cyclic_eta_block_decomposition}. Neither the source formula
nor this theorem asserts independence, a Gram relation, kernel stabilization,
or rectangular vanishing for the closed tensors. Taking traces of cyclic
products returns the established positive-power trace identities.

Consequently, the strongest explicit common-matrix conclusion remains
\begin{align}
    T^2
        &= T^3,
    \label{eq:cpgsv17_audit_square_cube}\\
    \tr(T^N)
        &= \tr(T)
        \qquad (N\geq 1),
    \label{eq:cpgsv17_audit_constant_traces}
\end{align}
together with entrywise nonnegativity and primitivity. The formal
counterexample shows that these conclusions, even supplemented by injectivity,
full virtual spanning, and positivity of every neighboring operator, do not
imply Eq.~\ref{eq:cpgsv17_audit_required_vanishing}. The explicit Hayashi
decomposition and inverse-map calculation show that the same nonzero nilpotent
remainder occurs for the source-selected factorization. The all-cut Markov
decomposition proves the strong area law for the same raw tensor. This tensor
satisfies the displayed SAL and literal ZCL relations, but it is not the
injective normal tensor assumed in Case~I. Its normal representative loses
literal ZCL. Thus the example isolates the normalization-sensitive step
required for $T_{k,h}=a_kb_h$; it does not refute Lemma~C.5 under all standing
hypotheses.

This obstruction does not survive the full normal Case-I hypotheses: the
completed route below excludes it and formalizes the structural corollary at
source lines 1503--1506. It remains relevant to the later Case-II passage. The
sectorwise analytic part of that passage is complete. For every absorbed BNT
representative $\mathcal K_s=\mu_sA_s$, the formalization proves one-site
injectivity, positivity on every nonempty periodic chain, saturation of the
area law, and the literal identity
\begin{align}
    \mathcal T_{\mathcal K_s}^2
        &= \mathcal T_{\mathcal K_s}.
    \label{eq:cpgsv17_absorbed_literal_transfer_idempotence}
\end{align}
The remaining obstruction is normality after coefficient absorption. If
$A_s$ is normal, then the transfer map of $\mu_sA_s$ has spectral radius
$|\mu_s|^2$. The global canonical-form convention supplies at least one
unit-modulus coefficient, but it does not imply $|\mu_s|=1$ for every sector.
Consequently, the normal Case-I structural theorem cannot be applied to all
absorbed BNT representatives. The constructions from an explicitly supplied
rank-one neighboring trace factorization remain valid restricted results, but
they do not prove that such a factorization exists for every absorbed
representative.

The explicit example discussed in the conclusion realizes this normalization
obstruction under every condition used in the simple Case-II passage, at the
normalized fixed-representative boundary. It refutes the absorption inference,
not the source assertion that some suitable factorization exists.

\section{The unit-closure-trace necessary condition}
\label{sec:cpgsv17_unit_closure_trace}

The existential assertion of condition~(iv) quantifies over all possible
physical-sector factorizations of each absorbed representative, so a
counterexample must exclude every factorization and a proof may choose any.
The formalization records a constraint that every qualifying factorization
imposes on the tensor itself, independently of that choice.

Let $\mathcal K$ be any MPO tensor admitting a physical-sector factorization
with positive semidefinite neighboring operators $\eta_{k,h}$ and real
families satisfying
\begin{align}
    \tr(\eta_{k,h})
        &= a_kb_h,
    \label{eq:cpgsv17_neighboring_rank_one_trace}\\
    \sum_k a_kb_k
        &= 1.
    \label{eq:cpgsv17_neighboring_trace_normalization}
\end{align}
Closing the physical legs of the factorized slices and using unitary
invariance of trace writes the physical-trace transfer as a rectangular
product
\begin{align}
    \mathcal T_{\mathcal K}
        &= LQ,
    \label{eq:cpgsv17_physical_trace_rectangular_product}\\
    L_{\beta,k}
        &:= \tr((l_k)_\beta),
    \label{eq:cpgsv17_left_trace_matrix}\\
    Q_{k,\alpha}
        &:= \tr((r_k)_\alpha).
    \label{eq:cpgsv17_right_trace_matrix}
\end{align}
The opposite product is the rank-one coefficient matrix:
\begin{align}
    (QL)_{k,h}
        &= \tr(\eta_{k,h})
    \label{eq:cpgsv17_opposite_product_neighboring_trace}\\
        &= a_kb_h.
    \label{eq:cpgsv17_opposite_product_rank_one}
\end{align}
Cyclic invariance of trace gives
\begin{align}
    \tr((LQ)^N)
        &= \tr((QL)^N)
        \qquad (N\geq 1).
    \label{eq:cpgsv17_rectangular_power_trace_cyclicity}
\end{align}
All positive powers of the normalized rank-one matrix equal the matrix itself,
which has trace one. Write $\rho^{(N)}(\mathcal K)$ for the periodic operator
denoted $\sigma^{(N)}(\mathcal K)$ in the source. Hence
\begin{align}
    \tr(\mathcal T_{\mathcal K}^{\,N})
        &= 1
        \qquad (N\geq 1),
    \label{eq:cpgsv17_unit_transfer_power_trace}\\
    \tr\rho^{(N)}(\mathcal K)
        &= 1
        \qquad (N\geq 1).
    \label{eq:cpgsv17_unit_periodic_operator_trace}
\end{align}
The formal statements are
\leanid{MPOTensor.PhysicalSectorFactorization.physTraceTransfer_eq_leftTraceMatrix_mul_rightTraceMatrix},
\leanid{MPOTensor.PhysicalSectorFactorization.NeighboringTraceFactorization.trace_physTraceTransfer_pow_eq_one},
and
\leanid{MPOTensor.PhysicalSectorFactorization.NeighboringTraceFactorization.trace_mpo_eq_one}
in \path{TNLean/MPS/MPDO/NeighboringTraceObstruction.lean}.

Three consequences follow. First, the normalization $\tr(T)=1$ invoked at
source line 1498 of Ref.~\cite{Cirac2016MPDO_arXiv} is a constraint on the
tensor rather than a convention: a tensor with a closed chain of non-unit
trace admits no factorization of the required form. Second, the
per-representative assertion of condition~(iv) carries the implicit
normalization $\tr(\mathcal T_{\mathcal K_s})=1$ for each absorbed
representative, whereas the ambient normalization at source lines 1755--1759
constrains only the weighted sum over representatives. Whether the standing
hypotheses of condition~(ii) imply this per-representative unit trace is part
of the open content of the all-sector structure statement. Any admissible
tensor violating it immediately refutes the printed per-representative
assertion. Third, under literal zero correlation length, the condition says
that the physical-trace transfer of each absorbed representative is an
idempotent of unit trace. The four-sector active-trace witness is consistent
with this condition because its physical-trace transfer is a rank-one
coordinate projection.

\section{Conclusion}
\label{sec:cpgsv17_conclusion}

The proof in Appendix C.2 of Ref.~\cite{Cirac2016MPDO_arXiv} contains an
unjustified matrix step as written: a primitive entrywise nonnegative matrix
with constant traces of positive powers need not have rank one. The obstruction
is a nilpotent Jordan component at the zero eigenvalue, which the trace
identities do not detect.

The formal development proves two corrected sufficient matrix theorems. The
first derives the rank-one factorization from positive semidefiniteness, trace
normalization, and the trace-power identities. The second derives $T^2=T$ from
the sector-pairing identity and linear independence of the closed sector
tensors; trace one then gives the same rank-one factorization directly. These
criteria remain available as independent matrix results. The unrelated
aggregate bundle that paired a three-site Markov witness with an independently
supplied trace matrix has been retired: it imposed no relation among the
state, the matrix, and an MPO tensor and therefore added no consequence to
either criterion.

The completed normal Case-I route supplies a third, tensor-attached result at
the source hypothesis boundary. The active-sector conclusion is formalized in
\leanid{MPOTensor.exists_activeSectorTraceMatrix_rank_one_coefficients_of_isSAL_of_literal_ZCL},
and the all-sector coefficient statement is formalized in
\leanid{MPOTensor.exists_neighboringOperator_trace_rank_one_coefficients_of_isSAL_of_literal_ZCL}.
Normality and literal ZCL, together with the physical-sector positivity and
primitivity data, force the active sector to be unique and hence $T^2=T$. For
the project factors $T=QL$, this gives
\begin{align}
    Q(\Id-LQ)L
        &= QL-(QL)^2
    \label{eq:cpgsv17_normal_case_remainder_difference}\\
        &= 0.
    \label{eq:cpgsv17_normal_case_remainder_vanishing}
\end{align}
Thus the nilpotent-zero obstruction is excluded on this normal Case-I surface;
idempotence of $LQ$ alone would still yield only $T^2=T^3$.

For the all-sector statement, retain the active coefficient functions on
\begin{align}
    I_*
        &:= \{k:p_k\neq 0\}
    \label{eq:cpgsv17_active_index_set}
\end{align}
and set
\begin{align}
    a_k
        &= 0
        \qquad (p_k=0),
    \label{eq:cpgsv17_inactive_left_coefficient}\\
    b_k
        &= 0
        \qquad (p_k=0).
    \label{eq:cpgsv17_inactive_right_coefficient}
\end{align}
In the zero-weight-reparameterized and rephased factorization selected by the
construction, $p_k=0$ makes both every incoming and every outgoing neighboring
operator vanish. On active pairs, the all-sector theorem first gives
\begin{align}
    \operatorname{Re}\tr(\eta_{k,h})
        &= a_kb_h.
    \label{eq:cpgsv17_active_real_neighboring_trace}
\end{align}
Every neighboring operator is positive semidefinite and hence Hermitian with
real trace. Therefore,
\begin{align}
    \tr(\eta_{k,h})
        &= a_kb_h
    \label{eq:cpgsv17_all_sector_complex_trace_factorization}
\end{align}
for every pair of sector labels, and filtering the diagonal sum to $I_*$ gives
\begin{align}
    \sum_k a_kb_k
        &= 1.
    \label{eq:cpgsv17_all_sector_coefficient_normalization}
\end{align}
This is a project-derived realization of the printed all-index coefficient
equation in Lemma~C.5 of Ref.~\cite{Cirac2016MPDO_arXiv}: CPSV16 quantifies
sector labels without defining TNLean's active--inactive split or its ambient
filler sectors.

The Case-II sectorwise analytic slice is formalized separately. Literal
physical-trace idempotence gives copy independence in the source standing
simple-canonical form through
\leanid{MPOTensor.weight_copy_independent_of_isPhysicalTraceIdempotent}. It
first descends to every weighted canonical block through
\leanid{MPOTensor.weighted_basis_physTraceTransfer_sq_of_literal_ZCL}, and then
to every absorbed BNT representative through
\leanid{MPOTensor.commonWeightAbsorbedBasisMPOTensor_physTraceTransfer_sq_of_literal_ZCL}.
Together with the projector, positivity, and entropy arguments,
\leanid{MPOTensor.commonWeightAbsorbedBasisMPOTensor_caseII_properties_of_literal_ZCL}
gives, for every $s$,
\begin{align}
    \mathcal K_s
        &\ \text{injective},
    \label{eq:cpgsv17_caseii_absorbed_injective}\\
    \sigma^{(N)}(\mathcal K_s)
        &\geq 0
        \qquad (N\geq 1),
    \label{eq:cpgsv17_caseii_periodic_positive}\\
    \mathcal K_s
        &\ \text{satisfies SAL},
    \label{eq:cpgsv17_caseii_absorbed_sal}\\
    \mathcal T_{\mathcal K_s}^2
        &= \mathcal T_{\mathcal K_s}.
    \label{eq:cpgsv17_caseii_literal_transfer_idempotence}
\end{align}
This literal identity is distinct from the scale-invariant predicate
\leanid{MPOTensor.IsSourceZCL}. The proof temporarily derives that predicate
with scalar $1$ only to invoke the existing entropy theorem; it retains the
literal identity as a separate conclusion.

The absorbed tensor
\begin{align}
    \mathcal K_s
        &= \mu_sA_s
    \label{eq:cpgsv17_absorbed_bnt_representative}
\end{align}
is an absorbed BNT representative, not automatically a normal representative.
Normality of $A_s$ fixes the spectral radius of its transfer map to one,
whereas coefficient absorption changes that radius to $|\mu_s|^2$. The
source's global normalization gives only one unit-modulus coefficient across
the whole canonical form and does not supply $|\mu_s|=1$ for every $s$.

The printed absorbed-normality step is formally refuted under all conditions
used in Appendix C.2 of Ref.~\cite{Cirac2016MPDO_arXiv} by the two-sector,
bond-one theorem
\leanid{MPOTensor.CaseIIAbsorptionCounterexample.%
ambient_isSAL_isSimpleCanonicalForm_and_firstAbsorbed_not_isNormalTensor}.
For every $N>0$, its diagonal periodic operator is positive semidefinite and
has trace two. Explicit trace-preserving completely positive maps witness the
Definition~4.1 closure equations, hence SAL, and its physical-trace transfer
is the identity, giving literal ZCL. Its normal representatives have disjoint
doubled-physical support, global weights $(1/\sqrt2,1)$, normalized
fixed-representative simple canonical form, and biCF. Nevertheless, the first
absorbed representative has transfer spectral radius $1/2$ and is not normal.
This repairs the missing full-condition statement and decides the printed
normality inference used in Proposition~\texttt{prop2to3}; it does not refute
the asserted all-sector conclusion or the literal implication
(iv)$\Rightarrow$(v).

That implication is separately refuted by the one-representative, two-copy
construction in
\path{TNLean/MPS/MPDO/Theorem49RepeatedCopyCounterexample.lean},
\leanid{MPOTensor.CaseIIAbsorptionCounterexample.%
printed_theorem49_iv_to_v_is_false}. Its raw copy weights are $1$ and $1/2$.
The theorem packages the exact BNT canonical assembly and global unit-weight
condition, MPDO positivity, nonnilpotence of the sole BNT representative,
simultaneous one-letter span, representative MPDO positivity, the
distinct-layer equation, the physical-sector factorization, and the positive
rank-one neighboring-trace data. Thus every printed standing hypothesis and
all of condition~(iv) have explicit component witnesses. The two-site block
also generates MPDOs. Nevertheless, it fails Definition~4.1 of
Ref.~\cite{Cirac2016MPDO_arXiv}: transfer idempotence would equate
\begin{align}
    1+2^{-4}
        &= 1+2^{-2},
    \label{eq:cpgsv17_repeated_copy_false_trace_equality}
\end{align}
which is false. The literal implication is therefore an \emph{unfaithful
statement}; treating outer-BNT assembly as its only unfinished step is
\emph{proof-path drift}.

Proposition~\texttt{prop2to5} of Ref.~\cite{Cirac2016MPDO_arXiv} remains a
viable missing statement because it starts from condition~(ii). The earlier
Case-II argument at source lines 1646--1665 uses literal ZCL to make
repeated-copy weights independent of the copy label and absorb their common
value before the short proof of \texttt{prop2to5} invokes Proposition~C.7. For
the rescaling-stable example, the separate explicit vertical identification
is complete in
\leanid{MPOTensor.RescalingStableLengthDependentRFP.R_oneLabelBNTAlgebraTensorClause};
it does not supply sectorwise normality and is not used in this descent
argument.

The structural corollary at source lines 1503--1506 is formalized by
\leanid{MPOTensor.exists_physicalSectorFactorization_rank_one_coefficients_of_isSAL_of_literal_ZCL}.
It retains the selected physical-sector factorization itself: an isometry $U$
with TNLean's orientation
\begin{align}
    U\kappa_{\beta,\alpha}U^\dagger
        &= \bigoplus_k
        (l_k)_\beta\otimes(r_k)_\alpha,
    \label{eq:cpgsv17_physical_sector_orientation}
\end{align}
its left and right factor matrices, and ambient positive semidefinite
neighboring operators. Together with
Eqs.~\ref{eq:cpgsv17_all_sector_complex_trace_factorization} and
\ref{eq:cpgsv17_all_sector_coefficient_normalization}, this completes the
normal Case-I corollary on the literal-ZCL surface. It asserts no additional
closure property of the sector tensors. The orientation in
Eq.~\ref{eq:cpgsv17_physical_sector_orientation} is the convention of
\leanid{MPOTensor.PhysicalSectorFactorization}, rather than an identification
with a separately drawn graphical orientation in the source.

A second boundary specialization is complete without normality. For virtual
bond dimension one,
\leanid{MPOTensor.exists_physicalSectorFactorization_of_bondDim_one}
constructs one physical sector with left factor $\mathbb C^d$ and
one-dimensional right factor. Its neighboring operator is the one-site
periodic operator and is therefore positive semidefinite by SAL. The
physical-trace transfer is a nonzero scalar by SAL and equals one by literal
idempotence, so the neighboring trace has the normalized factorization
\begin{align}
    1
        &= 1\cdot 1.
    \label{eq:cpgsv17_bond_one_trace_factorization}
\end{align}
This is a conditional helper at the bond-one boundary. It does not derive the
factorization for absorbed representatives of arbitrary virtual dimension, so
\ghissue{6775} and its unrestricted Blueprint node remain open.

This one-factorization result is a \emph{conditional helper}. It assumes a
neighboring trace factorization of the whole tensor to which the channels are
applied. Condition~(iv) instead supplies a separate physical factorization for
each BNT representative and leaves its raw repeated-copy weights
unconstrained. The scalar two-copy example shows that no direct-sum
construction can turn these hypotheses alone into conclusion~(v).

The projector-controlled construction remains relevant to
Proposition~\texttt{prop2to5}, hence to the viable source implication
(ii)$\Rightarrow$(v). After ZCL has forced common copy weights and those
values have been absorbed, Proposition~C.7 constructs a channel pair within
each outer BNT factorization
\cite[Appendix~C.2]{Cirac2016MPDO_arXiv}. A direct concatenation of the inner
sector labels is not valid: cross-BNT neighboring operators vanish by support
orthogonality, while a single global rank-one equation
$\tr(\eta_{k,h})=a_kb_h$ would generally make cross terms nonzero. The
construction instead uses the outer support projectors to control a direct sum
of channel pairs. This all-sector assembly is formalized in \ghissue{6632},
separately from the refuted raw (iv)$\Rightarrow$(v) reading.

The arbitrary-boundary part of this construction is formalized by
\leanid{MPSTensor.trace_evalWord_gauge_toTensorFromBlocks_mul} and
\leanid{MPOTensor.%
physCloseN_eq_sum_commonWeightAbsorbedBasis_of_gauge_toTensor}. Let $G$ be the
invertible virtual matrix implementing the gauge equivalence between the raw
tensor and its sector decomposition, and let $X$ be an arbitrary virtual
boundary. The boundary entering a fixed representative is the sum of the
diagonal copy blocks of $G^{-1}XG$ over that representative. The same
transformation is valid at every chain length. Copy-weight equality is used
only among copies of one representative, so this calculation imposes no
relation between distinct BNT coefficients. The projector-controlled
combination is formalized by
\leanid{MPOTensor.%
blockTwo_isRFPViaTS_of_commonWeightAbsorbedBNTChannels}. For a complete outer
projective resolution, the measured projection on one blocked site is the
reindexed matrix $P_s\otimes\Id$. Compression selects the matching absorbed
representative at both required closure lengths; its supplied channel is then
applied.

The completeness hypothesis is removed by
\leanid{MPOTensor.%
blockTwo_isRFPViaTS_of_commonWeightAbsorbedBNTChannels_of_orthogonalSupports}.
For pairwise orthogonal two-sided supports,
\begin{align}
    Q
        &:= \sum_sP_s
    \label{eq:cpgsv17_outer_support_projection}
\end{align}
is an orthogonal projection, and both relevant closures lie in its first-site
corner. The projector-controlled active maps are completely positive and
preserve the trace selected by $Q$. Measuring $\Id-Q$ and preparing a fixed
density matrix therefore completes each active map to a channel, while the
complementary term vanishes on the closures. This resolves \ghissue{6807}
without asserting that the active Markov projection is the ambient identity.
Thus the outer-sector channel combination is complete. The source proof now
has one remaining obligation: the source-context factorization in
\ghissue{6775}. The direct sector-mixing problem formerly separated as
\ghissue{6793} is a project-derived alternative, not another step asserted or
used by Appendix C.2 of Ref.~\cite{Cirac2016MPDO_arXiv}.

The composition from supplied factorization data is complete end to end.
Given, for every absorbed representative, a physical-sector factorization
with positive semidefinite neighboring operators and normalized rank-one trace
coefficients, the theorem
\leanid{MPOTensor.%
blockTwo_isRFPViaTS_of_bntLayerOrthogonal_of_physicalSectorFactorization%
_of_commonWeight} constructs the pairwise orthogonal outer supports from
layer orthogonality and one-site spanning, produces each representative
channel pair through Proposition~C.7 of
Ref.~\cite{Cirac2016MPDO_arXiv}, and assembles them into the blocked channel
pair of the complete tensor. The coarsest merging is packaged generically by
\leanid{MPOTensor.PhysicalSectorFactorization.NeighboringTraceFactorization.%
ofOneSector}: a single-sector factorization whose sole neighboring operator is
positive semidefinite with unit trace suffices. The sole remaining obligation
of the source implication (ii)$\Rightarrow$(v) is therefore the existence of
these per-representative factorizations.

The all-cut Markov proof and the source-selected inverse-map calculation show
that the raw four-sector tensor satisfies the displayed SAL, literal ZCL, and
inverse-map relations while the trace matrix of those four fixed sectors is
not rank one. This is not an existential obstruction. Reindexing a physical
label as two binary labels gives a one-sector factorization with
\begin{align}
    l_0
        &= \tfrac14\Id,
    \label{eq:cpgsv17_one_sector_left_identity}\\
    l_1
        &= Z,
    \label{eq:cpgsv17_one_sector_left_pauli}\\
    r_0
        &= \Id,
    \label{eq:cpgsv17_one_sector_right_identity}\\
    r_1
        &= \tfrac18Z.
    \label{eq:cpgsv17_one_sector_right_pauli}
\end{align}
Its sole neighboring operator is
\begin{align}
    \eta
        &= \Id\otimes\tfrac14\Id
           +\tfrac18Z\otimes Z.
    \label{eq:cpgsv17_one_sector_neighboring_operator}
\end{align}
This operator is positive semidefinite and has trace one. The factorization
and its normalized trace coefficients are formalized in
\leanid{MPOTensor.ActiveSectorSpanningCounterexample.%
exists_oneSector_neighboringTraceFactorization}. The non-rank-one trace matrix
of the four selected inverse-map sectors therefore refutes only the proof path
tied to those sectors.

A different four-letter tensor refutes the corresponding low-level
implication. Define
\begin{align}
    L
        &:=
        \begin{pmatrix}
            1&1&1&1\\
            1&2&-7&4
        \end{pmatrix},
    \label{eq:cpgsv17_noncartesian_left_matrix}\\
    R
        &:=
        \begin{pmatrix}
            1/4&-3/100\\
            1/4&-1/100\\
            1/4&1/100\\
            1/4&3/100
        \end{pmatrix}.
    \label{eq:cpgsv17_noncartesian_right_matrix}
\end{align}
Take the four diagonal physical slices to be the outer products of the
corresponding columns and rows. Then
\begin{align}
    LR
        &= \operatorname{diag}(1,0),
    \label{eq:cpgsv17_noncartesian_physical_projection}
\end{align}
so the physical-trace transfer is literally idempotent, while $RL$ has rank
two. The four slices span $M_2(\mathbb C)$, and the displayed scalar sector
factorization has positive neighboring operators. Hence the tensor is
injective and satisfies SAL. Perron normalization expresses it as
\begin{align}
    \mathcal K
        &= \mu A,
    \label{eq:cpgsv17_noncartesian_perron_normalization}\\
    0
        &< |\mu|
         < 1,
    \label{eq:cpgsv17_noncartesian_subunit_weight}
\end{align}
where $A$ is normal.

The obstruction applies to every physical-sector factorization, not merely to
the displayed one. The scalar slice and the two slices with simple spectra
$1,2,-7,4$ and $-3,-1,1,3$ force every left and right sector factor to be
one-dimensional. The physical isometry then assigns each physical label to a
unique sector. The neighboring trace matrix consequently contains a nonzero
two-by-two minor, whereas a factorization
$\tr(\eta_{k,h})=a_kb_h$ has rank at most one. This is formalized in
\leanid{MPOTensor.NonCartesianActiveSectorCandidate.%
full_lowLevel_counterexample}. Thus the local properties already inherited by
an absorbed representative do not suffice to prove the factorization.
Precisely, the formal theorem supplies injectivity, SAL and hence MPDO
generation, literal physical-trace idempotence, and an expression
$\mathcal K=\mu A$ with $A$ normal and $0<|\mu|<1$. It does not package an
ambient simple tensor in biCF whose canonical reconstruction contains
$\mathcal K$.

For the singleton normalization produced here,
\begin{align}
    |\mu|
        &= \sqrt r,
    \label{eq:cpgsv17_noncartesian_weight_perron_value}
\end{align}
where $r$ is the Perron value of the unnormalized transfer. For the positive
test matrix
\begin{align}
    H
        &:= \operatorname{diag}(20,1),
    \label{eq:cpgsv17_positive_test_matrix}
\end{align}
direct calculation of the adjoint action gives
\begin{align}
    H-\mathcal E_{\mathcal K}^*(H)
        &=
        \begin{pmatrix}
            85/8&-9/40\\
            -9/40&4697/5000
        \end{pmatrix}
        >0.
    \label{eq:cpgsv17_adjoint_transfer_positive_gap}
\end{align}
Its $(1,1)$ entry is positive, and its determinant is
\begin{align}
    \det(H-\mathcal E_{\mathcal K}^*(H))
        &= \frac{85}{8}\frac{4697}{5000}
           -\left(\frac9{40}\right)^2
    \label{eq:cpgsv17_positive_gap_determinant_expansion}\\
        &= \frac{19861}{2000}
         >0,
    \label{eq:cpgsv17_positive_gap_determinant}
\end{align}
which certifies positive definiteness. Pairing
Eq.~\ref{eq:cpgsv17_adjoint_transfer_positive_gap} with a positive-definite
Perron eigenmatrix $\rho$ of $\mathcal E_{\mathcal K}$ gives
\begin{align}
    0
        &<
        \tr\bigl(        (H-\mathcal E_{\mathcal K}^*(H))\rho
        \bigr)
    \label{eq:cpgsv17_positive_gap_pairing}\\
        &=
        (1-r)\tr(H\rho).
    \label{eq:cpgsv17_trace_adjoint_perron_identity}
\end{align}
Since $\tr(H\rho)>0$, Eq.~\ref{eq:cpgsv17_trace_adjoint_perron_identity}
gives $r<1$ and hence $|\mu|<1$. This is formalized in
\leanid{MPOTensor.NonCartesianActiveSectorCandidate.%
transferMap_posDef_eigenvalue_lt_one}. Since this normalization has only one
coefficient, the source convention at line 246 requires that coefficient to
have unit modulus and therefore fails. Rescaling to unit weight destroys
literal ZCL. This is the scale tension recorded in
\path{docs/paper-gaps/cpsv16_unit_weight_rfp_scale_tension.tex}. Consequently,
the construction does not resolve \ghissue{6775} or refute the printed
source-context assertion.

The non-Cartesian tensor is not itself the normal tensor assumed in Case~I,
and its normal representative loses literal ZCL. It therefore does not refute
Lemma~C.5 under its complete Case-I hypotheses. The completed normal Case-I
corollary does not replace literal ZCL by TNLean's scale-invariant
\leanid{MPOTensor.IsSourceZCL} predicate. The completed Case-II analytic
inheritance likewise does not prove that coefficient absorption preserves
normality.

By contrast, the earlier four-sector active-trace witness $\mathcal K$ has
explicit trace-preserving completely positive renormalization maps, both
before and after two-site blocking, but its sole BNT coefficient is strictly
subunit. Its globally normalized representative
$A=(4/\sqrt5)\mathcal K$ retains SAL and TNLean's scale-invariant
\leanid{MPOTensor.IsSourceZCL} relation, while $A^{[2]}$ fails the literal
idempotence forced by conclusion~(v). The normalized representative therefore
refutes the broadened implication using that up-to-scalar predicate, but it is
not an unqualified counterexample to Theorem~4.9 of
Ref.~\cite{Cirac2016MPDO_arXiv} because it fails the paper's literal ZCL
diagram.

Independently, the repeated-copy theorem above refutes the exact printed
implication (iv)$\Rightarrow$(v). The source-context
(ii)$\Rightarrow$(iv) factorization remains \ghissue{6775}. Once it is
supplied, the Proposition~\texttt{3to5} construction and its twice-applied
channels for one representative are formalized by
\leanid{MPOTensor.PhysicalSectorFactorization.NeighboringTraceFactorization.%
blockTwo_isRFPViaTS}. Combining those supplied channel pairs through the
orthogonal outer physical sectors then gives (ii)$\Rightarrow$(v); that
combination is formalized in \ghissue{6632}, and its trace-preserving
completion outside the active support resolves \ghissue{6807}. The
project-derived direct sector-mixing alternative in \ghissue{6793} is retired
in favor of this source route.

\bibliographystyle{alpha}
\bibliography{references}

\end{document}
